% Keep paragraph alignment independent of declarations in earlier chapters.
\begingroup
\leftskip=0pt
\rightskip=0pt
\parfillskip=0pt plus 1fil
\chapter{Komplexe Zahlen}
\label{chap:complex_numbers}

Komplexe Zahlen erweitern die reellen Zahlen. Ihre geometrische Darstellung
führt vom Lösen quadratischer Gleichungen zu Potenzen, Wurzeln und
allgemeinen Polynomgleichungen.

% Handschrift: Teil 7 - Komplexe Zahlen, Seiten 1--8.
\section{Von den reellen Zahlen zur komplexen Zahlenebene}
\label{sec:complex_numbers_cartesian_form}

Komplexe Zahlen vereinfachen viele Rechnungen und treten in der Beschreibung
physikalischer Zusammenhänge auf. Zunächst motiviert uns eine Gleichung,
deren Koeffizienten reell sind, die aber keine reelle Lösung besitzt:
\[
  z^2-6z+25\overset{!}{=}0
  \quad\Longleftrightarrow\quad (z-3)^2=-16.
\]
Auch die $pq$-Formel führt auf $3\pm\sqrt{9-25}$ und damit auf eine negative
Zahl unter der Wurzel. Können wir den Zahlbereich so erweitern, dass diese
und schließlich alle nichtkonstanten Polynomgleichungen lösbar werden?

\par\noindent\begin{minipage}{\linewidth}
\begin{DefinitionBox}{Imaginäre Einheit und komplexe Zahlen}{complex_numbers_cartesian}
  Die \emph{imaginäre Einheit} \EnglishTerm{imaginary unit} $i$ erfüllt
  $i^2=-1$. Da keine reelle Zahl diese Eigenschaft hat, gilt $i\notin\R$.

  Eine \emph{komplexe Zahl} \EnglishTerm{complex number} besitzt die
  \emph{kartesische Darstellung} \EnglishTerm{Cartesian form}
  \[
    z=x+iy,\qquad x,y\in\R,\qquad
    \C=\{x+iy\mid x,y\in\R\}.
  \]
  Dabei heißen $\operatorname{Re}(z)=x$ der \emph{Realteil}
  \EnglishTerm{real part} und $\operatorname{Im}(z)=y$ der
  \emph{Imaginärteil} \EnglishTerm{imaginary part}.
\end{DefinitionBox}
\vspace{-0.5\baselineskip}
\begin{ExampleBox}{Realteil und Imaginärteil ablesen}{complex_numbers_components}
  Bestimmen Sie die beiden reellen Komponenten der folgenden Zahlen.
  \begin{center}
    \renewcommand{\arraystretch}{1.3}
    \begin{tabular}{@{}c@{\qquad}c@{\qquad}c@{}}
      $z$ & $\operatorname{Re}(z)$ & $\operatorname{Im}(z)$\\\midrule
      $2+3i$ & $\Cloze{2cm}{2}$ & $\Cloze{2cm}{3}$\\
      $i$ & $\Cloze{2cm}{0}$ & $\Cloze{2cm}{1}$\\
      $3$ & $\Cloze{2cm}{3}$ & $\Cloze{2cm}{0}$
    \end{tabular}
  \end{center}
\end{ExampleBox}
\end{minipage}\par


\begin{ExampleBox}{Die Ausgangsgleichung lösen}{complex_numbers_initial_equation}
  Bestimmen Sie mithilfe von $i^2=-1$ die beiden Lösungen der Ausgangsgleichung.
  \[
    (z-3)^2=(4i)^2
      \quad\Longleftrightarrow\quad z=\Cloze{3cm}{3\pm4i}.
  \]
  Die Probe bestätigt $(\pm4i)^2=16i^2=-16$.
\end{ExampleBox}

\clearpage
\begin{RemarkBox}[remark title=Realteil und Imaginärteil sind reell]
  Realteil und Imaginärteil sind beide \emph{reell}; der Imaginärteil von
  $2+3i$ ist $3$, nicht $3i$.
  Zahlen mit Imaginärteil $0$ sind reell, Zahlen mit Realteil $0$ heißen
  \emph{rein imaginär} \EnglishTerm{purely imaginary}. Damit gilt
  \[
    z\in\R\ \Longleftrightarrow\ \operatorname{Im}(z)=0,
    \qquad \N\subsetneq\Z\subsetneq\Q\subsetneq\R\subsetneq\C.
  \]
\end{RemarkBox}

\subsection{Zahlenebene und Betrag}
Die \emph{Gaußsche Zahlenebene} \EnglishTerm{complex plane} ordnet
$z=x+iy$ dem Punkt $(x,y)$ zu. Die waagerechte \emph{reelle Achse}
\EnglishTerm{real axis} trägt den Realteil, die senkrechte
\emph{imaginäre Achse} \EnglishTerm{imaginary axis} den Imaginärteil.
Jede reelle Zahl liegt auf der reellen Achse, jede rein imaginäre Zahl
auf der imaginären Achse.
\begin{center}
  \CachedTikzFigure{content/resources/chapters/complex_numbers/tikz/complex_plane_cartesian.tex}
\end{center}
Der Pfeil vom Ursprung zum Punkt veranschaulicht den Abstand vom Ursprung.
Aus dem rechtwinkligen Dreieck folgt mit dem Satz des Pythagoras die Betragsformel.

\par\noindent\begin{minipage}{\linewidth}
\begin{DefinitionBox}{Betrag einer komplexen Zahl}{complex_numbers_modulus}
  Der \emph{Betrag} \EnglishTerm{modulus} einer komplexen Zahl $z=x+iy$
  ist ihr Abstand vom Ursprung:
  \[
    |z|=\sqrt{x^2+y^2}
       =\sqrt{\bigl(\operatorname{Re}(z)\bigr)^2+
                    \bigl(\operatorname{Im}(z)\bigr)^2}\in\R_{\ge0}.
  \]
\end{DefinitionBox}
\vspace{-0.5\baselineskip}
\begin{ExampleBox}{Eine Zahl als Punkt und als Pfeil}{complex_numbers_plane_example}
  Bestimmen Sie den Betrag von $z=3+2i$ anhand der eingezeichneten Komponenten.
  \par\smallskip
  \begin{minipage}[c]{0.50\linewidth}
    \CachedTikzFigure{content/resources/chapters/complex_numbers/tikz/complex_plane_modulus.tex}
  \end{minipage}\hfill
  \begin{minipage}[c]{0.45\linewidth}
    Die Katheten haben die Längen $3$ und $2$. Der Abstand beträgt daher
    \[
      |3+2i|=\Cloze{3cm}{\sqrt{3^2+2^2}=\sqrt{13}}.
    \]
  \end{minipage}
\end{ExampleBox}
\end{minipage}\par

\clearpage
\subsection{Komplexe Konjugation}
Zur Konjugation betrachten wir die Spiegelung eines Punktes an der reellen Achse.

\par\noindent\begin{minipage}{\linewidth}
\begin{DefinitionBox}{Konjugiert komplexe Zahl}{complex_numbers_conjugate}
  Zu $z=x+iy$ heißt $z^*=x-iy$ die \emph{konjugiert komplexe Zahl}
  \EnglishTerm{complex conjugate}. Sie entsteht durch
  \emph{Spiegelung an der reellen Achse}.
  \par\medskip
  \begin{minipage}[c]{0.51\linewidth}
    \CachedTikzFigure{content/resources/chapters/complex_numbers/tikz/complex_conjugation.tex}
  \end{minipage}\hfill
  \begin{minipage}[c]{0.44\linewidth}
    \textbf{Einfach jedes $i$ durch $-i$ ersetzen.}

    Der Realteil bleibt unverändert; der Imaginärteil wechselt sein Vorzeichen:
    \[
      \begin{aligned}
        \operatorname{Re}(z^*)&=\operatorname{Re}(z),\\[4pt]
        \operatorname{Im}(z^*)&=-\operatorname{Im}(z).
      \end{aligned}
    \]
    Diese Regel gilt auch in der später eingeführten Exponentialdarstellung.
  \end{minipage}
\end{DefinitionBox}
\vspace{-0.5\baselineskip}
\begin{ExampleBox}{Konjugation und Betrag}{complex_numbers_conjugation_example}
  Bestimmen Sie zu $z=3+2i$ die konjugierte Zahl und vergleichen Sie die Beträge.
  \[
    z^*=\Cloze{3cm}{3-2i},\qquad
    |z|=\sqrt{3^2+2^2}=\sqrt{13},\qquad
    |z^*|=\Cloze{3cm}{\sqrt{3^2+(-2)^2}=\sqrt{13}}.
  \]
\end{ExampleBox}
\end{minipage}\par

\begin{KeyPointBox}{Eigenschaften der Konjugation}{complex_numbers_conjugation_properties}
  Zweimaliges Spiegeln führt zurück zur Ausgangszahl; Punkte auf der
  Spiegelachse bleiben fest. Daher gilt
  \[
    (z^*)^*=z,\qquad
    z=z^*\ \Longleftrightarrow\ \operatorname{Im}(z)=0
          \ \Longleftrightarrow\ z\in\R.
  \]
  Die Spiegelung ändert den Abstand vom Ursprung nicht:
  \[
    |z^*|=\sqrt{x^2+(-y)^2}=\sqrt{x^2+y^2}=|z|.
  \]
\end{KeyPointBox}

\clearpage
% Handschrift: Seiten 8--13; Betragsregeln von Seite 6 hier eingeordnet.
\section{Rechnen in kartesischer Darstellung}
\label{sec:complex_numbers_cartesian_arithmetic}
Wir rechnen mit den vertrauten algebraischen Gesetzen und ersetzen dabei
$i^2$ durch $-1$. Addition und Subtraktion lassen sich unmittelbar in der
Zahlenebene deuten; die Geometrie der Multiplikation folgt später.

\begin{KeyPointBox}{Addition und Subtraktion}{complex_numbers_addition}
  Für $z_1=x_1+iy_1$ und $z_2=x_2+iy_2$ werden die beiden Anteile
  getrennt addiert beziehungsweise subtrahiert:
  \[
    z_1\pm z_2=(x_1\pm x_2)+i(y_1\pm y_2).
  \]
  Geometrisch entspricht dies der Addition beziehungsweise Subtraktion
  der zugehörigen Pfeile.
\end{KeyPointBox}
\begin{ExampleBox}{Summe und Differenz in der Zahlenebene}{complex_numbers_sum_difference}
  Bestimmen Sie für $z_1=2+i$ und $z_2=3-2i$ die Summe und die Differenz.
  \[
    \begin{aligned}
      z_1+z_2&=(2+3)+i(1-2)=\Cloze{3cm}{5-i},\\[5pt]
      z_1-z_2&=(2-3)+i(1+2)=\Cloze{3cm}{-1+3i}.
    \end{aligned}
  \]
  \begin{center}
    \CachedTikzFigure{content/resources/chapters/complex_numbers/tikz/complex_addition_subtraction.tex}
  \end{center}
\end{ExampleBox}
\clearpage
Für eine Zahl und ihre konjugierte Zahl heben sich beim Addieren die
Imaginärteile und beim Subtrahieren die Realteile auf:
\[
  \begin{aligned}
    z+z^*&=(x+iy)+(x-iy)=2x
      &&\Longrightarrow& \operatorname{Re}(z)&=\frac{z+z^*}{2},\\[6pt]
    z-z^*&=(x+iy)-(x-iy)=2iy
      &&\Longrightarrow& \operatorname{Im}(z)&=\frac{z-z^*}{2i}.
  \end{aligned}
\]
\begin{ExampleBox}{Komponenten durch Konjugation zurückgewinnen}{complex_numbers_components_conjugation}
  Prüfen Sie beide Formeln an $z=3+2i$ und $z^*=3-2i$.
  \[
    \operatorname{Re}(z)=\frac{(3+2i)+(3-2i)}2=\Cloze{1.5cm}{3},\qquad
    \operatorname{Im}(z)=\frac{(3+2i)-(3-2i)}{2i}=\Cloze{1.5cm}{2}.
  \]
\end{ExampleBox}

\subsection{Multiplikation}
Ausmultiplizieren und Zusammenfassen nach Real- und Imaginärteil liefert
\[
  \begin{aligned}
    z_1z_2&=(x_1+iy_1)(x_2+iy_2)\\[5pt]
      &=x_1x_2+ix_1y_2+iy_1x_2+i^2y_1y_2\\[5pt]
      &=(x_1x_2-y_1y_2)+i(x_1y_2+y_1x_2).
  \end{aligned}
\]
\begin{RemarkBox}[remark title=Gemischte Terme bei der Multiplikation]
  Real- und Imaginärteil eines Produkts sind im Allgemeinen \emph{nicht}
  die Produkte der einzelnen Real- und Imaginärteile. Die gemischten
  Terme und $i^2=-1$ müssen berücksichtigt werden.
\end{RemarkBox}
\begin{KeyPointBox}{Produkt mit der konjugierten Zahl}{complex_numbers_conjugate_product}
  Beim Multiplizieren einer Zahl mit ihrer konjugierten Zahl verschwinden
  die gemischten Terme:
  \[
    zz^*=(x+iy)(x-iy)=x^2-i^2y^2=x^2+y^2=|z|^2\in\R_{\ge0}.
  \]
  Damit lässt sich der Betrag auch durch $|z|=\sqrt{zz^*}$ ausdrücken.
\end{KeyPointBox}
\begin{ExampleBox}{Den Betrag durch ein Produkt prüfen}{complex_numbers_modulus_product_example}
  Berechnen Sie für $z=3+2i$ das Produkt $zz^*$ und vergleichen Sie
  mit der geometrisch bestimmten Länge $\sqrt{13}$.
  \begin{ClozeArea}{4}
    \[
      (3+2i)(3-2i)=9-(2i)^2=9+4=13,
      \qquad |z|=\sqrt{13}.
    \]
  \end{ClozeArea}
\end{ExampleBox}

\clearpage
\subsection{Division}
Um einen Quotienten kartesisch darzustellen, machen wir seinen Nenner
mithilfe der komplexen Konjugation reell.
\begin{KeyPointBox}{Erweitern mit dem konjugierten Nenner}{complex_numbers_division}
  Für $z_2\ne0$ gilt mit $z_2z_2^*=|z_2|^2$ die Umformung
  \[
    \begin{aligned}
      \frac{z_1}{z_2}
        &=\frac{z_1z_2^*}{|z_2|^2}
         =\frac{(x_1+iy_1)(x_2-iy_2)}{x_2^2+y_2^2}\\[7pt]
        &=\frac{x_1x_2+y_1y_2}{x_2^2+y_2^2}
          +i\,\frac{y_1x_2-x_1y_2}{x_2^2+y_2^2}.
    \end{aligned}
  \]
  Auch beim Dividieren dürfen Real- und Imaginärteile nicht getrennt
  durcheinander geteilt werden.
\end{KeyPointBox}

\begin{ExampleBox}{Einen Quotienten kartesisch darstellen}{complex_numbers_quotient_example}
  Bringen Sie $(3+2i)/(1-i)$ auf die Form $x+iy$, indem Sie den Nenner reell machen.
  \begin{ClozeArea}{6}
    \[
      \begin{aligned}
        \frac{3+2i}{1-i}
          &=\frac{(3+2i)(1+i)}{(1-i)(1+i)}\\[6pt]
          &=\frac{3+5i+2i^2}{1-i^2}
           =\frac{1+5i}{2}=\frac12+\frac52i.
      \end{aligned}
    \]
  \end{ClozeArea}
  Realteil und Imaginärteil des Quotienten sind also $1/2$ und $5/2$.
\end{ExampleBox}
\begin{KeyPointBox}{Die imaginäre Einheit im Nenner}{complex_numbers_reciprocal_i}
  Beim Wechsel von $i$ aus dem Nenner in den Zähler entsteht ein Minuszeichen:
  \[
    \frac1i=\frac{-i}{i(-i)}=\frac{-i}{1}=-i,
    \qquad \frac{a}{i}=-ia\quad(a\in\C).
  \]
\end{KeyPointBox}

\clearpage
\subsection{Rechengesetze}
Die aus dem Reellen bekannten Gesetze gelten auch für komplexe Zahlen:
\[
  \begin{array}{@{}ll@{}}
    z_1+z_2=z_2+z_1,\quad z_1z_2=z_2z_1
      &\text{Kommutativgesetze},\\[7pt]
    z_1+(z_2+z_3)=(z_1+z_2)+z_3
      &\text{Assoziativgesetze},\\[7pt]
    z_1(z_2z_3)=(z_1z_2)z_3&\\[7pt]
    z_1(z_2+z_3)=z_1z_2+z_1z_3
      &\text{Distributivgesetz}.
  \end{array}
\]
Für die Konjugation ergeben sich außerdem die Beziehungen
\[
  (z_1+z_2)^*=z_1^*+z_2^*,\quad
  (z_1z_2)^*=z_1^*z_2^*,\quad
  \left(\frac{z_1}{z_2}\right)^*=\frac{z_1^*}{z_2^*}\quad(z_2\ne0).
\]
\begin{ExampleBox}{Die Reihenfolge des Zusammenfassens wählen}{complex_numbers_associativity}
  Berechnen Sie $(1+i)(3+2i)(3-2i)$ mit zwei verschiedenen Klammerungen.
  Nutzen Sie im ersten Weg das konjugierte Paar.
  \begin{ClozeArea}{7}
    \[
      \begin{aligned}
        (1+i)\bigl((3+2i)(3-2i)\bigr)
          &=(1+i)\cdot13=13+13i,\\[6pt]
        \bigl((1+i)(3+2i)\bigr)(3-2i)
          &=(1+5i)(3-2i)\\[6pt]
          &=3-2i+15i-10i^2=13+13i.
      \end{aligned}
    \]
  \end{ClozeArea}
\end{ExampleBox}

\begin{KeyPointBox}{Rechenregeln für den Betrag}{complex_numbers_modulus_laws}
  Für komplexe Zahlen gelten die Produkt- und Quotientenregel sowie
  die \emph{Dreiecksungleichung} \EnglishTerm{triangle inequality}:
  \[
    |z|\ge0,\qquad |z_1z_2|=|z_1||z_2|,\qquad
    \left|\frac{z_1}{z_2}\right|=\frac{|z_1|}{|z_2|}\quad(z_2\ne0),
  \]
  \[
    |z_1+z_2|\le|z_1|+|z_2|.
  \]
\end{KeyPointBox}

\clearpage
% Handschrift: Seiten 14--18; Argument geometrisch statt Falltabelle (Autorwunsch).
\section{Polardarstellung und Darstellungswechsel}
\label{sec:complex_numbers_polar_form}
Statt durch zwei Koordinaten lässt sich eine komplexe Zahl durch ihren
Abstand vom Ursprung und ihre Richtung beschreiben. Diese Sicht bereitet
die geometrische Deutung von Produkten und Potenzen vor.

\begin{DefinitionBox}{Trigonometrische Polardarstellung}{complex_numbers_polar}
  Für $z\ne0$ heißen $r=|z|>0$ und ein Winkel $\varphi$ zur positiven
  reellen Achse \emph{Polarkoordinaten} \EnglishTerm{polar coordinates}.
  Der Winkel heißt \emph{Argument} \EnglishTerm{argument} oder Phasenwinkel.
  Die \emph{Polardarstellung} \EnglishTerm{polar form} lautet
  \[
    z=r(\cos\varphi+i\sin\varphi).
  \]
  Für $z=0$ ist kein Argument festgelegt.
  \begin{center}
    \CachedTikzFigure{content/resources/chapters/complex_numbers/tikz/complex_polar_coordinates.tex}
  \end{center}
\end{DefinitionBox}
Die rechtwinklige Zerlegung verbindet Polarkoordinaten und kartesische Koordinaten:
\[
  \cos\varphi=\frac{x}{r},\quad \sin\varphi=\frac{y}{r}
  \quad\Longrightarrow\quad
  x=r\cos\varphi,\quad y=r\sin\varphi.
\]
\begin{ExampleBox}{Betrag und Winkel in der Zahlenebene}{complex_numbers_polar_examples}
  Zeichnen Sie die Zahlen $z_1=3(\cos45^\circ+i\sin45^\circ)$ und
  $z_2=4(\cos150^\circ+i\sin150^\circ)$ ausgehend von Betrag und Winkel.
  \begin{center}
    \CachedTikzFigure{content/resources/chapters/complex_numbers/tikz/complex_polar_examples.tex}
  \end{center}
  Die Pfeillängen betragen $3$ beziehungsweise $4$; die Winkel werden
  von der positiven reellen Achse aus gemessen.
\end{ExampleBox}

\clearpage
\subsection{Volle Umdrehungen und besondere Winkel}
Eine volle Umdrehung verändert die Zahl nicht. Deshalb beschreiben
$\varphi$ und $\varphi+2k\pi$ mit $k\in\Z$ dieselbe Richtung:
\[
  \cos40^\circ+i\sin40^\circ
  =\cos400^\circ+i\sin400^\circ.
\]
Häufig wählt man einen Winkel in $[0,2\pi)$ beziehungsweise $[0^\circ,360^\circ)$.
Auch negative Winkel sind möglich; sie werden im Uhrzeigersinn gemessen.
\begin{ExampleBox}{Die vier Achsenrichtungen}{complex_numbers_trigonometric_axes}
  Lesen Sie die Werte aus der Lage auf dem Einheitskreis ab.
  \begin{center}
    \renewcommand{\arraystretch}{1.5}
    \begin{tabular}{@{}c@{\qquad}c@{\qquad}c@{}}
      $\varphi$ & $\cos\varphi+i\sin\varphi$ & Lage\\\midrule
      $0^\circ$ & $\Cloze{2cm}{1}$ & positive reelle Achse\\
      $90^\circ$ & $\Cloze{2cm}{i}$ & positive imaginäre Achse\\
      $180^\circ$ & $\Cloze{2cm}{-1}$ & negative reelle Achse\\
      $270^\circ$ & $\Cloze{2cm}{-i}$ & negative imaginäre Achse\\
      $360^\circ$ & $\Cloze{2cm}{1}$ & wieder die Ausgangsrichtung
    \end{tabular}
  \end{center}
\end{ExampleBox}

\subsection{Von der Polarform zur kartesischen Form}
Zur Umrechnung berechnen wir die beiden Komponenten $x=r\cos\varphi$
und $y=r\sin\varphi$ getrennt.
\begin{ExampleBox}{Komponenten aus Betrag und Winkel bestimmen}{complex_numbers_polar_to_cartesian}
  Bestimmen Sie Real- und Imaginärteil der Zahl
  $z=3(\cos45^\circ+i\sin45^\circ)$ aus dem vorigen Beispiel.
  \begin{ClozeArea}{5}
    \[
      x=3\cos45^\circ=\frac3{\sqrt2},\qquad
      y=3\sin45^\circ=\frac3{\sqrt2},\qquad
      z=\frac3{\sqrt2}(1+i).
    \]
  \end{ClozeArea}
\end{ExampleBox}

\clearpage
\subsection{Von der kartesischen Form zum Argument}
Der Betrag ist unmittelbar $r=\sqrt{x^2+y^2}$. Für den Winkel nutzen wir
bei $x\ne0$ die Beziehung
\[
  \tan\varphi=\frac{\sin\varphi}{\cos\varphi}
    =\frac{y/r}{x/r}=\frac{y}{x}.
\]
\begin{KeyPointBox}{Der Arkustangens kennt den Quadranten nicht}{complex_numbers_argument_geometry}
  $\arctan(y/x)$ liefert nur einen Winkel zwischen $-\pi/2$ und $\pi/2$.
  Gegenüberliegende Pfeile haben denselben Quotienten $y/x$ und werden
  dadurch nicht unterschieden.

  \textbf{Stellen Sie sich $z$ in der Zahlenebene vor oder zeichnen Sie es ein.}
  Prüfen Sie, ob der berechnete Winkel in die richtige Richtung zeigt.
  Zeigt er genau entgegengesetzt, addieren oder subtrahieren Sie
  $\pi$ beziehungsweise $180^\circ$ passend zur Skizze.

  Bei $x=0$ ist $y/x$ nicht definiert. Lesen Sie den Winkel direkt an der
  imaginären Achse ab: nach oben $\pi/2$, nach unten $-\pi/2$
  (gleichwertig $3\pi/2$). Für $z=0$ gibt es kein Argument.
\end{KeyPointBox}
\begin{ExampleBox}{Ein Punkt im dritten Quadranten}{complex_numbers_argument_example}
  Bestimmen Sie Betrag und einen passenden Winkel von $z=-1-i$.
  Vergleichen Sie die tatsächliche Richtung mit dem Ergebnis des Arkustangens.
  \par\smallskip
  \begin{minipage}[c]{0.47\linewidth}
    \CachedTikzFigure{content/resources/chapters/complex_numbers/tikz/complex_argument_quadrant.tex}
  \end{minipage}\hfill
  \begin{minipage}[c]{0.48\linewidth}
    Der Punkt liegt links unten. Dagegen zeigt
    $\arctan((-1)/(-1))=\pi/4$ nach rechts oben.
    Wir drehen diesen Pfeil um eine halbe Umdrehung.
    \begin{ClozeArea}[15]{6}
      \[
        r=\sqrt2,\qquad
        \varphi=\frac\pi4-\pi=-\frac{3\pi}{4}.
      \]
      Mit $+\pi$ ergibt sich gleichwertig $5\pi/4$.
    \end{ClozeArea}
  \end{minipage}
  \par\medskip
  Die trigonometrische Darstellung lautet damit
  \[
    z=\sqrt2\left(\cos\left(-\frac{3\pi}{4}\right)
                +i\sin\left(-\frac{3\pi}{4}\right)\right)
     =\sqrt2\left(\cos\frac{3\pi}{4}-i\sin\frac{3\pi}{4}\right).
  \]
\end{ExampleBox}

\clearpage
% Handschrift: Seiten 19--26.
\section{Euler-Formel und Exponentialdarstellung}
\label{sec:complex_numbers_exponential_form}
Die Euler-Formel verbindet Trigonometrie und Exponentialfunktion.
Mit ihr werden aus Winkeladditionen einfache Potenzgesetze.

\begin{TheoremBox}{Euler-Formel}{complex_numbers_euler}
  Für jeden reellen Winkel $\varphi$ im Bogenmaß gilt die
  \emph{Euler-Formel} \EnglishTerm{Euler's formula}
  \[
    e^{i\varphi}=\cos\varphi+i\sin\varphi.
  \]
\end{TheoremBox}
Der Beweis mithilfe von Potenzreihen folgt in Mathematik~2.
Konjugieren bedeutet auch hier, jedes $i$ durch $-i$ zu ersetzen:
\[
  (e^{i\varphi})^*=e^{-i\varphi}
    =\cos\varphi-i\sin\varphi.
\]
Über die Formeln für Real- und Imaginärteil erhalten wir daraus
\[
  \cos\varphi=\operatorname{Re}(e^{i\varphi})
    =\frac{e^{i\varphi}+e^{-i\varphi}}2,\qquad
  \sin\varphi=\operatorname{Im}(e^{i\varphi})
    =\frac{e^{i\varphi}-e^{-i\varphi}}{2i}.
\]
\begin{RemarkBox}[remark title=Vergleich mit Hyperbelfunktionen]
  Diese Darstellungen ähneln den Definitionen der
  \emph{Hyperbelfunktionen} \EnglishTerm{hyperbolic functions}:
  \[
    \begin{array}{@{}cc@{}}
      \text{Trigonometrisch}&\text{Hyperbolisch}\\[5pt]
      \displaystyle\sin x=\frac{e^{ix}-e^{-ix}}{2i}
        &\displaystyle\sinh x=\frac{e^x-e^{-x}}2\\[12pt]
      \displaystyle\cos x=\frac{e^{ix}+e^{-ix}}2
        &\displaystyle\cosh x=\frac{e^x+e^{-x}}2
    \end{array}
  \]
  Man spricht von Sinus hyperbolicus und Kosinus hyperbolicus.
\end{RemarkBox}
\begin{ExampleBox}{Eine trigonometrische Identität nachweisen}{complex_numbers_double_angle}
  Zeigen Sie mit der Exponentialdarstellung die Doppelwinkelformel
  $2\sin\varphi\cos\varphi=\sin(2\varphi)$.
  \begin{ClozeArea}{7}
    \[
      \begin{aligned}
        2\sin\varphi\cos\varphi
          &=\frac{(e^{i\varphi}-e^{-i\varphi})(e^{i\varphi}+e^{-i\varphi})}{2i}\\[7pt]
          &=\frac{e^{2i\varphi}-e^{-2i\varphi}}{2i}
           =\sin(2\varphi).
      \end{aligned}
    \]
  \end{ClozeArea}
\end{ExampleBox}

\clearpage
\subsection{Vom Einheitskreis zur Exponentialdarstellung}
Die Zahl $e^{i\varphi}$ besitzt den Betrag $1$. Dies lässt sich auf zwei
Wegen nachrechnen:
\[
  |e^{i\varphi}|=\sqrt{\cos^2\varphi+\sin^2\varphi}=1,
  \qquad
  |e^{i\varphi}|=\sqrt{e^{i\varphi}e^{-i\varphi}}=\sqrt{e^0}=1.
\]
Multiplikation mit $r>0$ streckt den Pfeil vom Einheitskreis auf die Länge $r$.
\begin{center}
  \CachedTikzFigure{content/resources/chapters/complex_numbers/tikz/complex_exponential_scaling.tex}
\end{center}
\begin{DefinitionBox}{Exponentialdarstellung}{complex_numbers_exponential}
  Mit der Euler-Formel wird aus der trigonometrischen Polardarstellung die
  \emph{Exponentialdarstellung} \EnglishTerm{exponential form}:
  \[
    z=r(\cos\varphi+i\sin\varphi)=re^{i\varphi},\qquad r=|z|>0.
  \]
  Der Winkel $\varphi$ im Exponenten wird immer im \emph{Bogenmaß} angegeben.
  Für die Konjugation gilt $z^*=re^{-i\varphi}$.
\end{DefinitionBox}
\begin{ExampleBox}{Besondere Winkel}{complex_numbers_special_angles}
  Ordnen Sie die Winkel den Achsenrichtungen zu und ergänzen Sie die Werte.
  \[
    e^{i0}=\Cloze{1cm}{1},\quad
    e^{i\pi/2}=\Cloze{1cm}{i},\quad
    e^{i\pi}=\Cloze{1cm}{-1},\quad
    e^{3i\pi/2}=\Cloze{1cm}{-i},\quad
    e^{2\pi i}=\Cloze{1cm}{1}.
  \]
\end{ExampleBox}
Volle Umdrehungen lassen die Zahl auch in Exponentialdarstellung unverändert:
\[
  re^{i(\varphi+2k\pi)}
    =re^{i\varphi}(e^{2\pi i})^k=re^{i\varphi},\qquad k\in\Z.
\]
So beschreiben $2e^{-i\pi/4}$ und $2e^{7i\pi/4}$ dieselbe Zahl:
$-45^\circ$ und $315^\circ$ zeigen in dieselbe Richtung.

\clearpage
\subsection{Addition: zuerst die Darstellung wechseln}
Bei Summen und Differenzen rechnen wir zweckmäßig in kartesischer Form.
Beträge und Winkel dürfen dabei nicht getrennt addiert oder subtrahiert werden.
\begin{ExampleBox}{Eine Summe wieder in Exponentialform angeben}{complex_numbers_exponential_sum}
  Bestimmen Sie $z_1+z_2$ für $z_1=2e^{i\pi}$ und $z_2=e^{i\pi/2}$.
  Wechseln Sie zunächst in die kartesische Form und bestimmen Sie danach
  Betrag und Richtung des Ergebnisses.
  \[
    z_1=-2,\qquad z_2=i,\qquad z_1+z_2=-2+i.
  \]
  \par\smallskip
  \begin{minipage}[c]{0.46\linewidth}
    \CachedTikzFigure{content/resources/chapters/complex_numbers/tikz/complex_sum_argument.tex}
  \end{minipage}\hfill
  \begin{minipage}[c]{0.49\linewidth}
    Die Summe liegt links oben. $\arctan(-1/2)$ zeigt dagegen nach rechts
    unten. Eine halbe Umdrehung gegen den Uhrzeigersinn liefert den
    gesuchten Winkel.
    \begin{ClozeArea}[16]{7}
      \[
        \begin{aligned}
          r&=\sqrt{(-2)^2+1^2}=\sqrt5,\\[5pt]
          \varphi&=\arctan(-1/2)+\pi\\[5pt]
                 &\approx2{,}678\approx153{,}4^\circ.
        \end{aligned}
      \]
    \end{ClozeArea}
  \end{minipage}
  \par\medskip
  Somit ist $z_1+z_2=\sqrt5\,e^{i(\pi-\arctan(1/2))}$.
\end{ExampleBox}

\subsection{Multiplikation und Division}
Für $z_1=r_1e^{i\varphi_1}$ und $z_2=r_2e^{i\varphi_2}$ folgen aus den
Potenzgesetzen die beiden Rechenregeln
\[
  z_1z_2=r_1r_2e^{i(\varphi_1+\varphi_2)},\qquad
  \frac{z_1}{z_2}=\frac{r_1}{r_2}e^{i(\varphi_1-\varphi_2)}\quad(z_2\ne0).
\]
\begin{KeyPointBox}{Geometrische Bedeutung}{complex_numbers_rotation}
  Beim Multiplizieren werden die Beträge multipliziert und die Winkel addiert.
  Beim Dividieren werden die Beträge dividiert und die Winkel subtrahiert.
  Multiplikation mit $i=e^{i\pi/2}$ dreht um $90^\circ$ gegen den Uhrzeigersinn;
  Division durch $i$ dreht um $90^\circ$ im Uhrzeigersinn.
\end{KeyPointBox}

\clearpage
\begin{ExampleBox}{Produkt und Quotient geometrisch verstehen}{complex_numbers_product_example}
  Für dieselben Zahlen $z_1=2e^{i\pi}=-2$ und $z_2=e^{i\pi/2}=i$
  bestimmen wir nun Produkt und Quotient. Die Länge bleibt jeweils $2$.
  \begin{ClozeArea}{6}
    \[
      \begin{aligned}
        z_1z_2&=2e^{i\pi}e^{i\pi/2}
                =2e^{3i\pi/2}=-2i,\\[8pt]
        \frac{z_1}{z_2}&=\frac{2e^{i\pi}}{e^{i\pi/2}}
                =2e^{i(\pi-\pi/2)}=2e^{i\pi/2}=2i.
      \end{aligned}
    \]
  \end{ClozeArea}
  \begin{center}
    \CachedTikzFigure{content/resources/chapters/complex_numbers/tikz/complex_multiplication_rotation.tex}
  \end{center}
  Links dreht sich der Pfeil von $z_1$ um $+\pi/2$, rechts um $-\pi/2$.
\end{ExampleBox}

\clearpage
% Handschrift: Seiten 27--33.
\section{Potenzen und komplexe Wurzeln}
\label{sec:complex_numbers_powers_and_roots}
Wiederholtes Multiplizieren potenziert den Betrag und vervielfacht den
Winkel. Beim umgekehrten Problem müssen wir berücksichtigen, dass Winkel
nur bis auf volle Umdrehungen festgelegt sind.

\begin{KeyPointBox}{Ganzzahlige Potenzen}{complex_numbers_powers}
  Für $z=re^{i\varphi}\ne0$ gilt
  \[
    z^n=\underbrace{z\cdots z}_{n\text{ Faktoren}}
       =r^ne^{in\varphi}\quad(n\ge1),\qquad
    z^0=1,\qquad z^{-n}=\frac1{z^n}.
  \]
  Der Betrag wird potenziert, der Winkel mit $n$ multipliziert.
  Für $z=0$ sind die positiven ganzzahligen Potenzen ebenfalls definiert.
\end{KeyPointBox}
Für $z=e^{i\varphi}$ bleibt der Betrag bei jeder Potenz gleich $1$.
Die Winkel $\varphi,2\varphi,3\varphi,\ldots$ folgen in gleichen Schritten:
\begin{center}
  \CachedTikzFigure{content/resources/chapters/complex_numbers/tikz/complex_unit_circle_powers.tex}
\end{center}
\begin{ExampleBox}{Eine hohe Potenz geschickt berechnen}{complex_numbers_eighth_power}
  Berechnen Sie $(1+i)^8$ in Exponentialform und vergleichen Sie mit
  wiederholtem Quadrieren. Der Punkt $1+i$ liegt rechts oben auf der
  Winkelhalbierenden, also ist $\varphi=\pi/4$ ohne weitere Korrektur.
  \begin{ClozeArea}{7}
    \[
      r=\sqrt{1^2+1^2}=\sqrt2,\qquad 1+i=\sqrt2\,e^{i\pi/4}.
    \]
    Daher gilt $(1+i)^8=(\sqrt2)^8e^{2\pi i}=16$.
    In kartesischer Form führt wiederholtes Quadrieren ebenfalls zum Ziel:
    \[
      (1+i)^2=2i,\qquad (1+i)^4=-4,\qquad (1+i)^8=16.
    \]
  \end{ClozeArea}
\end{ExampleBox}

\clearpage
Die Zwischenpotenzen zeigen zugleich die wachsenden Beträge und die Drehung.
Für $z=1+i$ liegen $z^2$, $z^4$, $z^6$ und $z^8$ auf den Koordinatenachsen:
\begin{center}
  \CachedTikzFigure{content/resources/chapters/complex_numbers/tikz/complex_numeric_powers.tex}
\end{center}
Eine volle Umdrehung ist bei $z^8$ erreicht; die Länge ist dann $16$.

\clearpage
\subsection{Komplexe Wurzeln als Lösungen einer Gleichung}
\begin{DefinitionBox}{Komplexe Wurzeln}{complex_numbers_roots}
  Für $a\in\C$ und $n\in\N$, $n\ge1$, heißt jede Lösung von $z^n=a$
  eine \emph{$n$-te Wurzel} \EnglishTerm{$n$th root} von $a$.
\end{DefinitionBox}
\vspace{-0.5\baselineskip}
\begin{ExampleBox}{Zwei Quadratwurzeln von $i$}{complex_numbers_square_roots_i}
  Finden Sie zwei Zahlen, deren Quadrat $i=e^{i\pi/2}$ ist.
  Beim Quadrieren wird der Winkel verdoppelt.
  \begin{ClozeArea}{5}
    \[
      z_0=e^{i\pi/4},\qquad z_1=e^{5i\pi/4},\qquad
      z_0^2=e^{i\pi/2}=i,\qquad z_1^2=e^{5i\pi/2}=i.
    \]
  \end{ClozeArea}
  \begin{center}
    \CachedTikzFigure{content/resources/chapters/complex_numbers/tikz/complex_square_roots_i.tex}
  \end{center}
  Die Lösungen liegen einander auf dem Einheitskreis gegenüber:
  Ihr Winkelabstand ist $\Delta\varphi=2\pi/2=\pi=180^\circ$.
\end{ExampleBox}

\begin{ExampleBox}{Vier vierte Wurzeln von $1$}{complex_numbers_fourth_roots_unity}
  Bestimmen Sie vier Zahlen mit $z^4=1$ und prüfen Sie durch Potenzieren.
  \par\smallskip
  \begin{minipage}[c]{0.48\linewidth}
    \CachedTikzFigure{content/resources/chapters/complex_numbers/tikz/complex_fourth_roots_unity.tex}
  \end{minipage}\hfill
  \begin{minipage}[c]{0.47\linewidth}
    \begin{ClozeArea}[15]{8}
      \[
        \begin{aligned}
          z_0&=1, &z_0^4&=1,\\[6pt]
          z_1&=e^{i\pi/2}=i, &z_1^4&=e^{2\pi i}=1,\\[6pt]
          z_2&=e^{i\pi}=-1, &z_2^4&=e^{4\pi i}=1,\\[6pt]
          z_3&=e^{3i\pi/2}=-i, &z_3^4&=e^{6\pi i}=1.
        \end{aligned}
      \]
    \end{ClozeArea}
  \end{minipage}
  \par\medskip
  Alle Lösungen liegen auf dem Einheitskreis, mit Winkelabstand
  $\Delta\varphi=2\pi/4=\pi/2=90^\circ$.
\end{ExampleBox}

\clearpage
\subsection{Alle Wurzeln systematisch bestimmen}
Schreiben wir $a=Re^{i\alpha}\ne0$ und $z=\rho e^{i\varphi}$ mit
$R,\rho>0$, so liefert der Vergleich von Betrag und Winkel
\[
  z^n=a
  \quad\Longleftrightarrow\quad
  \rho^n=R\ \text{ und }\ n\varphi=\alpha+2k\pi
  \text{ für ein }k\in\Z.
\]
Daraus folgen $\rho=\sqrt[n]{R}$ und $\varphi=(\alpha+2k\pi)/n$.
Durch Einsetzen prüfen wir die so erhaltenen Zahlen:
\[
  \left(\sqrt[n]{R}\,e^{i(\alpha+2k\pi)/n}\right)^n
    =Re^{i\alpha}e^{2k\pi i}=a.
\]
\begin{TheoremBox}{Die $n$ komplexen Wurzeln}{complex_numbers_root_formula}
  Für $a=Re^{i\alpha}\ne0$ sind die $n$ verschiedenen Lösungen von $z^n=a$
  \[
    z_k=\sqrt[n]{R}\,e^{i(\alpha+2k\pi)/n},\qquad k=0,\ldots,n-1.
  \]
  Sie liegen gleichmäßig auf dem Kreis mit Radius $\sqrt[n]{R}$;
  benachbarte Lösungen haben den Winkelabstand
  $\Delta\varphi=2\pi/n=360^\circ/n$.
  Für $a=0$ gibt es nur die Lösung $z=0$.
\end{TheoremBox}
Weitere ganze Werte von $k$ liefern keine neuen Lösungen, denn
\[
  z_{k+n}=\sqrt[n]{R}\,e^{i(\alpha+2k\pi+2n\pi)/n}
    =z_k e^{2\pi i}=z_k.
\]

\begin{ExampleBox}{Dritte Wurzeln von $-\sqrt3+i$}{complex_numbers_cube_roots}
  Bestimmen Sie alle Lösungen von $z^3=a=-\sqrt3+i$ und deuten Sie ihre Lage.
  Zunächst bringen wir die rechte Seite in Exponentialform.

  Der Punkt $a$ liegt links oben. Der Winkel
  $\arctan(1/(-\sqrt3))=-\pi/6$ zeigt jedoch nach rechts unten.
  Daher addieren wir eine halbe Umdrehung, also $\pi$.
  \begin{ClozeArea}{8}
    \[
      R=\sqrt{3+1}=2,\qquad
      \alpha=-\frac\pi6+\pi=\frac{5\pi}{6},\qquad
      a=2e^{5\pi i/6}.
    \]
    Die drei Wurzelwinkel sind $5\pi/18$, $17\pi/18$ und $29\pi/18$:
    \[
      z_0=\sqrt[3]2\,e^{5\pi i/18},\quad
      z_1=\sqrt[3]2\,e^{17\pi i/18},\quad
      z_2=\sqrt[3]2\,e^{29\pi i/18}.
    \]
  \end{ClozeArea}
  \begin{center}
    \CachedTikzFigure{content/resources/chapters/complex_numbers/tikz/complex_cube_roots.tex}
  \end{center}
  Die Lösungen liegen auf dem Kreis mit Radius $\sqrt[3]2$ im Winkelabstand
  $2\pi/3=120^\circ$. Beim Kubieren wird jeder ihrer Winkel auf
  $5\pi/6$ zuzüglich voller Umdrehungen abgebildet.
\end{ExampleBox}

\clearpage
% Handschrift: Seiten 34--36.
\section{Polynomgleichungen über den komplexen Zahlen}
\label{sec:complex_numbers_polynomial_equations}
Die Ausgangsfrage reicht über quadratische Gleichungen hinaus: Reicht die
Erweiterung auf $\C$ aus, um jede nichtkonstante Polynomgleichung zu lösen?
Der Fundamentalsatz der Algebra beantwortet diese Frage.

\begin{TheoremBox}{Fundamentalsatz der Algebra}{complex_numbers_fundamental_theorem}
  Jedes Polynom $p(z)=a_nz^n+\dots+a_1z+a_0$ mit komplexen Koeffizienten,
  $a_n\ne0$ und $n\ge1$, besitzt über $\C$ genau $n$ Nullstellen,
  \emph{mit ihren Vielfachheiten gezählt}. Es lässt sich vollständig zerlegen:
  \[
    p(z)=a_n(z-z_1)\cdots(z-z_n).
  \]
\end{TheoremBox}
Über $\R$ besitzt ein Polynom vom Grad $n$ höchstens $n$ Nullstellen.
Über $\C$ dürfen verschiedene Faktoren dieselbe Nullstelle besitzen. Die Anzahl
\emph{verschiedener} Nullstellen kann daher kleiner als $n$ sein.

\begin{ExampleBox}{Ein Polynom vollständig in Linearfaktoren zerlegen}{complex_numbers_polynomial_example}
  Bestimmen Sie alle Nullstellen von $p(z)=z^3-3z^2+4z-2$.
  Durch Einsetzen erkennen wir $p(1)=0$. Die Polynomdivision durch $z-1$
  wird daher ohne Rest aufgehen:
  \[
    \begin{array}{@{}r@{\,}r@{\,}r@{\,}r@{\quad}l@{}}
      (z^3&-3z^2&+4z&-2)&:(z-1)=z^2-2z+2\\[5pt]
      \underline{-(z^3}&\underline{-z^2)}&&&\\[5pt]
      &-2z^2&+4z&-2&\\[5pt]
      &\underline{-(-2z^2}&\underline{+2z)}&&\\[5pt]
      &&2z&-2&\\[5pt]
      &&\underline{-(2z}&\underline{-2)}&\\[5pt]
      &&&0&\text{Rest}
    \end{array}
  \]
  Somit gilt $p(z)=(z-1)(z^2-2z+2)$.
  Lösen Sie die verbleibende quadratische Gleichung und geben Sie die
  vollständige Linearfaktorzerlegung an.
  \begin{ClozeArea}{8}
    \[
      z^2-2z+2\overset{!}{=}0
      \quad\Longleftrightarrow\quad (z-1)^2=-1
      \quad\Longleftrightarrow\quad z=1\pm i.
    \]
    Somit sind $1$, $1+i$ und $1-i$ die drei Nullstellen und
    \[
      p(z)=(z-1)\bigl(z-(1+i)\bigr)\bigl(z-(1-i)\bigr).
    \]
  \end{ClozeArea}
\end{ExampleBox}

\par\noindent\begin{minipage}{\linewidth}
\begin{KeyPointBox}{Konjugierte Nullstellen bei reellen Koeffizienten}{complex_numbers_conjugate_roots}
  Hat $p$ nur reelle Koeffizienten, so gilt $p(z^*)=p(z)^*$.
  Mit jeder nichtreellen Nullstelle $z$ tritt deshalb auch $z^*$ als
  Nullstelle auf, jeweils mit derselben Vielfachheit.
\end{KeyPointBox}
Im Beispiel bilden $1+i$ und $1-i$ ein konjugiertes Paar. Ihr gemeinsamer
Faktor ist das reelle quadratische Polynom $z^2-2z+2$.

\begin{RemarkBox}[remark title=Geltungsbereich des Fundamentalsatzes]
  Der Fundamentalsatz gilt für Polynomgleichungen. Er besagt nicht,
  dass jede beliebige Gleichung in $\C$ lösbar ist: Die Gleichung
  $zz^*=-1$ hat keine Lösung, weil $zz^*=|z|^2\ge0$ gilt.
\end{RemarkBox}
\end{minipage}\par

\par\endgroup
